\documentclass[12pt, a4paper]{article}

\usepackage[utf8]{inputenc}
\usepackage[italian]{babel}
\usepackage{pgfplots}
\pgfplotsset{compat=1.18}
\usepackage{geometry}
\geometry{top=3.5cm, bottom=3.5cm, left=2.5cm, right=2.5cm}
\usepackage{xcolor}
% Evita righe fuori dai margini con i nomi lunghi in \texttt
\setlength{\emergencystretch}{3em}
\usepackage{listings}
\usepackage{float}
\usepackage{needspace}
\usetikzlibrary{arrows.meta, positioning}

% Stile delle figure con la pipeline dei job
\tikzset{
    pbox/.style={draw=green!50!black, line width=0.8pt, rounded corners=8pt,
        text width=6.6cm, minimum height=#1, align=left,
        font=\scriptsize\ttfamily, inner sep=6pt},
    pbox/.default=3.3cm,
    parrow/.style={-{Triangle[length=2.6mm, width=4mm]},
        line width=2.2pt, green!50!black}
}
% Titolo all'interno di ciascun riquadro della pipeline
\newcommand{\ptitle}[1]{{\normalfont\small\bfseries #1}\\[3pt]}
% Sei riquadri disposti a serpentina: 1 -> 2, giù, 4 <- 3, giù, 5 -> 6
\newcommand{\pipeline}[7]{%
\begin{tikzpicture}
\node[pbox=#1] (b1) at (0,0) {\ptitle{1. Input (file CSV)}#2};
\node[pbox=#1, right=1.4cm of b1] (b2) {\ptitle{2. Record letti (offset, riga)}#3};
\node[pbox=#1, below=1.2cm of b2] (b3) {\ptitle{3. Output del Map}#4};
\node[pbox=#1, below=1.2cm of b1] (b4) {\ptitle{4. Shuffle \& Sort}#5};
\node[pbox=#1, below=1.2cm of b4] (b5) {\ptitle{5. Reduce}#6};
\node[pbox=#1, right=1.4cm of b5] (b6) {\ptitle{6. Output finale}#7};
\draw[parrow] (b1) -- (b2);
\draw[parrow] (b2) -- (b3);
\draw[parrow] (b3) -- (b4);
\draw[parrow] (b4) -- (b5);
\draw[parrow] (b5) -- (b6);
\end{tikzpicture}}

% Stile del codice sorgente in appendice
\lstset{
    language=Java,
    basicstyle=\ttfamily\footnotesize,
    keywordstyle=\color{blue!70!black}\bfseries,
    commentstyle=\color{green!40!black}\itshape,
    stringstyle=\color{red!60!black},
    numbers=left,
    numberstyle=\tiny\color{gray},
    frame=single,
    breaklines=true,
    showstringspaces=false,
    tabsize=4,
    columns=fullflexible,
    keepspaces=true,
    literate={à}{{\`a}}1 {è}{{\`e}}1 {é}{{\'e}}1 {ì}{{\`i}}1 {ò}{{\`o}}1 {ù}{{\`u}}1
}

\begin{document}

\begin{titlepage}
    \centering
    
    % Intestazione
    {\scshape\LARGE Politecnico di Bari \par}
    \vspace{0.5cm}
    {\scshape\Large Ingegneria dei Sistemi Medicali \par}
    \vspace{1.5cm}
    {\Large Big Data Analytics \par}
    \vspace{2.5cm}
    
    % Titolo Progetto
    \rule{\linewidth}{0.5mm} \\[0.5cm]
    {\huge \bfseries Esonero Big Data \\[0.5cm] 
    \Large Analisi del Dataset YouTube in MapReduce \par}
    \rule{\linewidth}{0.5mm} \\[1.5cm]
    
    % Sottotitolo
    {\Large \textit{Report del Primo Progetto} \par}
    
    \vfill
    
    % Sezione inferiore: Docenti a sx, Studenti a dx
\noindent
\begin{minipage}[t]{0.40\textwidth}
    \vspace{0pt}
    \begin{flushleft}
        \large
        \emph{Docenti:}\\
        Prof. Simona \textsc{Colucci}\\
        Ing. Verdiana \textsc{Schena}
    \end{flushleft}
\end{minipage}
\hfill
\begin{minipage}[t]{0.55\textwidth}
    \vspace{0pt}
    \begin{flushright}
        \large
        \emph{Studenti:}\\
        Alessio \textsc{Troisi} - Matr. 601812\\
        Alice Maria \textsc{Patimo} - Matr. 603106\\
        Giovanni \textsc{Turi} - Matr. 602273
    \end{flushright}
\end{minipage}
    
    \vspace{2.5cm}
    
    %Anno Accademico
    {\large Novembre 2026 \par}
\end{titlepage}

% Indice generale del report
\tableofcontents
\newpage

\setcounter{section}{-1}
\section{Ambiente sperimentale e scelte generali}
\label{sec:generale}

In questa sezione sono descritte le scelte comuni ai tre job:
il dataset e i gruppi analizzati, l'ambiente di esecuzione,
l'organizzazione dei dati di input, la lettura dei file CSV
e il modo in cui sono stati raccolti log e risultati.
Le scelte specifiche di ciascun job sono invece discusse
nelle rispettive sezioni.

\subsection{Dataset e gruppi analizzati}
\label{sec:gruppi}

Il dataset comprende dieci Paesi: per ciascuno sono
disponibili un file CSV con i video in tendenza
e un file JSON con le categorie dei video
(Tabella~\ref{tab:file}).

\begin{table}[H]
\centering
\small
\begin{tabular}{|c|l|l|l|}
\hline
\textbf{Codice} & \textbf{Paese} & \textbf{File CSV} & \textbf{File JSON} \\
\hline
US & Stati Uniti & \texttt{USvideos.csv} & \texttt{US\_category\_id.json} \\
GB & Regno Unito & \texttt{GBvideos.csv} & \texttt{GB\_category\_id.json} \\
DE & Germania & \texttt{DEvideos.csv} & \texttt{DE\_category\_id.json} \\
CA & Canada & \texttt{CAvideos.csv} & \texttt{CA\_category\_id.json} \\
FR & Francia & \texttt{FRvideos.csv} & \texttt{FR\_category\_id.json} \\
RU & Russia & \texttt{RUvideos.csv} & \texttt{RU\_category\_id.json} \\
MX & Messico & \texttt{MXvideos.csv} & \texttt{MX\_category\_id.json} \\
KR & Corea del Sud & \texttt{KRvideos.csv} & \texttt{KR\_category\_id.json} \\
JP & Giappone & \texttt{JPvideos.csv} & \texttt{JP\_category\_id.json} \\
IN & India & \texttt{INvideos.csv} & \texttt{IN\_category\_id.json} \\
\hline
\end{tabular}
\caption{File del dataset.}
\label{tab:file}
\end{table}

Ogni record dei file CSV è composto dalle sedici colonne
riportate nella Tabella~\ref{tab:colonne}, nella quale
sono indicate anche le colonne utilizzate da ciascun job.

\begin{table}[H]
\centering
\small
\begin{tabular}{|c|l|l|}
\hline
\textbf{Indice} & \textbf{Colonna} & \textbf{Utilizzata da} \\
\hline
0 & \texttt{video\_id} & -- \\
1 & \texttt{trending\_date} & -- \\
2 & \texttt{title} & Job 1 \\
3 & \texttt{channel\_title} & Job 2 \\
4 & \texttt{category\_id} & Job 3 \\
5 & \texttt{publish\_time} & -- \\
6 & \texttt{tags} & -- \\
7 & \texttt{views} & Job 3 \\
8 & \texttt{likes} & Job 1 \\
9 & \texttt{dislikes} & Job 1 \\
10 & \texttt{comment\_count} & -- \\
11 & \texttt{thumbnail\_link} & -- \\
12 & \texttt{comments\_disabled} & -- \\
13 & \texttt{ratings\_disabled} & -- \\
14 & \texttt{video\_error\_or\_removed} & -- \\
15 & \texttt{description} & -- \\
\hline
\end{tabular}
\caption{Colonne dei file CSV.}
\label{tab:colonne}
\end{table}

Come richiesto dalla traccia, i job sono stati eseguiti
su tre gruppi di dataset di dimensione crescente:

\begin{itemize}
    \item \textbf{Dataset A} (3 Paesi): US, GB, DE.

    \item \textbf{Dataset B} (6 Paesi): US, GB, DE, CA, FR, RU.

    \item \textbf{Dataset C} (10 Paesi): US, GB, DE, CA, FR, RU,
    MX, KR, JP, IN.
\end{itemize}

\subsection{Ambiente di esecuzione}

Tutte le esecuzioni sono state effettuate su una
macchina virtuale Ubuntu con Hadoop 3.5.0.
Il codice, scritto in Java,
è stato esportato in un unico archivio \texttt{BigData.jar}
contenente le classi dei tre job.

Il progetto è gestito con Maven: le librerie di Hadoop
necessarie alla compilazione (\texttt{hadoop-common} e
\texttt{hadoop-mapreduce-client-core}, versione 3.5.0)
sono dichiarate nel file \texttt{pom.xml} e vengono
scaricate automaticamente, senza doverle importare a mano
su ogni macchina di sviluppo. Le dipendenze hanno scope
\texttt{provided}: servono soltanto in fase di compilazione,
perché durante l'esecuzione vengono fornite
dall'installazione di Hadoop; per questo motivo
non vengono incluse nel jar, che resta di pochi kilobyte.

Le tre configurazioni richieste dalla traccia
sono state ottenute attivando, di volta in volta,
le proprietà dei file di configurazione di Hadoop
(Tabella~\ref{tab:configurazioni}).

\begin{table}[H]
\centering
\footnotesize
\setlength{\tabcolsep}{4pt}
\begin{tabular}{|l|l|c|c|c|}
\hline
\textbf{File} & \textbf{Proprietà} & \textbf{Local} & \textbf{PD senza YARN} & \textbf{PD con YARN} \\
\hline
\texttt{core-site.xml} & \texttt{fs.defaultFS} & -- & HDFS & HDFS \\
\texttt{hdfs-site.xml} & \texttt{dfs.replication} & -- & 1 & 1 \\
\texttt{mapred-site.xml} & \texttt{mapreduce.framework.name} & -- & -- & \texttt{yarn} \\
\texttt{yarn-site.xml} & \texttt{aux-services} & -- & -- & \texttt{mapreduce\_shuffle} \\
\hline
\end{tabular}
\caption{Proprietà attive in ciascuna configurazione
(-- indica il valore di default).}
\label{tab:configurazioni}
\end{table}

\subsection{Organizzazione dei dati di input}

Per ciascun gruppo è stata creata un'unica cartella
(\texttt{dataset\_A}, \texttt{dataset\_B}, \texttt{dataset\_C})
contenente sia i file CSV sia i file JSON dei Paesi del gruppo,
utilizzata come input da tutti e tre i job.

Indicando una cartella come input, Hadoop elabora tutti
i file in essa contenuti, compresi i JSON, che servono
soltanto al Job 3. Le loro righe verrebbero scartate
dai controlli dei Mapper, perché contengono al massimo
due campi, ma ogni file JSON genererebbe un ulteriore
\textit{input split} e quindi un map task in più,
alterando i tempi di esecuzione e i contatori.
Per questo motivo i Driver scorrono il contenuto
della cartella e aggiungono come input soltanto
i file il cui nome termina con \texttt{videos.csv};
il Driver del Job 3 registra invece i file JSON
nella Distributed Cache (Sezione~\ref{sec:job3}).

Per le configurazioni pseudo-distribuite la stessa
struttura di cartelle è stata copiata su HDFS,
in modo che i comandi di esecuzione fossero identici
nelle tre configurazioni.

\subsection{Lettura dei record: la classe CSV\_Parser}
\label{sec:parser}

I campi testuali dei CSV, come i titoli, possono contenere
virgole. In questi casi il campo è racchiuso tra virgolette
e una semplice suddivisione sulla virgola produrrebbe
un numero errato di colonne.
La classe di supporto \texttt{CSV\_Parser}, condivisa
dai tre job, scorre la riga carattere per carattere
e considera separatori soltanto le virgole che si trovano
al di fuori delle virgolette.

Hadoop legge i file con \texttt{TextInputFormat},
che fornisce al Mapper una riga fisica alla volta.
Il campo \texttt{description}, ultima colonna dei CSV,
può però contenere ritorni a capo: un record logico
viene quindi letto come più righe.
Per questo motivo il numero di righe lette
(\textit{Map input records}) è superiore al numero
di record effettivi.
La prima riga di ciascun record contiene comunque
tutte le colonne che precedono la descrizione,
e quindi tutti i campi utilizzati dai job;
le righe successive (frammenti della descrizione)
non hanno la struttura attesa e vengono scartate
dal controllo sul numero di campi o dalla
conversione numerica.

\subsection{Struttura dei job}

Tutti i job seguono lo schema classico di MapReduce,
con tre classi: Mapper, Reducer e Driver.
Il Mapper emette una coppia chiave-valore per ciascun
record valido e non viene utilizzato un Combiner,
come confermano i contatori (\texttt{Combine input records=0}).
Gli oggetti \texttt{Text} e \texttt{Writable} emessi dai Mapper
vengono creati una sola volta e riutilizzati per ogni record,
evitando di allocare nuovi oggetti a ogni invocazione
del metodo \texttt{map()}.

L'aggregazione locale dei dati prima dello shuffle,
tramite un Combiner o direttamente nel Mapper,
è una possibile ottimizzazione: ridurrebbe il numero
di coppie trasferite al Reducer, in modo particolare
nei Job 2 e 3. Si è scelto di non adottarla per mantenere
lo schema visto durante il corso; il suo possibile impatto
è discusso nelle sezioni dei singoli job.

\subsection{Misura dei tempi di esecuzione}
\label{sec:tempi}

Ciascun Driver misura il tempo complessivo di esecuzione
con il metodo \texttt{System.\allowbreak currentTimeMillis()},
come differenza tra l'istante finale e quello iniziale,
espressa in millisecondi.
La misura comprende la configurazione, l'avvio
e l'esecuzione del job fino al suo completamento,
e viene stampata soltanto se il job termina con successo.

Ogni combinazione di job, configurazione e gruppo
di dataset è stata eseguita tre volte, per un totale
di 81 esecuzioni. Le tabelle riportano i tempi
delle singole esecuzioni e la loro media aritmetica,
arrotondata al millisecondo; i confronti e i grafici
si basano sulle medie.

Le esecuzioni sono state organizzate in tre giri:
in ciascun giro vengono eseguiti, nell'ordine, i tre job
sui tre gruppi di dataset, e soltanto al termine
del giro si passa alla ripetizione successiva.
In questo modo le tre misure di una stessa combinazione
non sono consecutive e risentono meno delle condizioni
momentanee della macchina.

Prima della serie in modalità Local i file dei dataset
sono stati letti una volta, in modo da caricarli
nella cache del sistema operativo.
Ciononostante, nella maggior parte dei casi la prima
esecuzione del Job 1, che apre ciascuna serie di misure,
è risultata la più lenta delle tre. Il rallentamento
riguarda quindi l'inizio della serie in sé più che
la lettura dei dati, ed è dovuto probabilmente
all'avvio dei primi processi Java e al caricamento
delle librerie di Hadoop (avvio \textit{a freddo}).

\subsection{Raccolta di log e risultati}

Inizialmente i log venivano copiati a mano dal terminale
e i risultati letti dai file \texttt{part-r-00000}
all'interno delle cartelle di output dei job.
Con 27 esecuzioni questo approccio è diventato scomodo
e soggetto a errori; sono stati quindi adottati
due accorgimenti.

\paragraph{Salvataggio dei log.}
L'output di ogni esecuzione è stato salvato con il comando
\texttt{tee}, che lo mostra nel terminale e allo stesso tempo
lo scrive in un file. Poiché Hadoop scrive i propri messaggi
(avanzamento e contatori) sullo \textit{standard error},
è necessario unirlo allo \textit{standard output}
con \texttt{2>\&1}:

\begin{verbatim}
hadoop jar BigData.jar com.BigData.Job3_Driver \
    datasets/dataset_A hadoop_out/job3-yarn-dataset_A \
    2>&1 | tee logs/job3-yarn-dataset_A.log
\end{verbatim}

\paragraph{Unione dei risultati.}
Per ottenere i risultati in un unico file è stato utilizzato
il comando \texttt{hadoop fs -getmerge}, che concatena
tutti i file \texttt{part-r-*} di una cartella di output
in un solo file locale:

\begin{verbatim}
hadoop fs -getmerge hadoop_out/job3-yarn-dataset_A \
    outputs/job3-yarn-dataset_A.txt
\end{verbatim}

Il comando funziona allo stesso modo sul filesystem locale
e su HDFS e ha permesso di trattare in modo uniforme
le tre configurazioni.
Accanto a ciascun risultato Hadoop crea un file nascosto
\texttt{.crc}, contenente le \textit{checksum} che il
filesystem di Hadoop utilizza per verificare l'integrità
dei dati; questi file sono stati conservati insieme
ai risultati.

Log e risultati seguono la stessa convenzione di nomi,
\texttt{job<N>-\allowbreak<configurazione>-\allowbreak<dataset>-\allowbreak run<K>},
con configurazione \texttt{local}, \texttt{pseudo} o \texttt{yarn}:

\begin{verbatim}
Esonero_Big_Data/
|-- BigData.jar
|-- datasets/      dataset_A/  dataset_B/  dataset_C/
|-- logs/          job1-local-dataset_A-run1.log  ...
|-- outputs/       job1-local-dataset_A-run1.txt  ...
`-- tempi/         tempi-local.csv  tempi-pseudo.csv  tempi-yarn.csv
\end{verbatim}

\subsection{Verifica della correttezza}

Per ciascun job e gruppo di dataset i file di output
prodotti nelle tre configurazioni sono risultati
identici byte per byte.

\section{Job 1 - Individuazione dei video con un numero di dislike maggiore del numero di like}
\subsection{Obiettivo}
Il primo job ha l'obiettivo di individuare i titoli dei video per i quali il numero di \textit{dislike} è maggiore del numero di \textit{like}.
L'analisi viene effettuata mediante il paradigma Hadoop MapReduce, attraverso tre componenti principali: Mapper, Reducer e Driver.
Il programma considera i campi \texttt{title}, \texttt{likes} e \texttt{dislikes} presenti nei file CSV dei video in tendenza (\textit{trending videos}) e seleziona i record in cui il numero di dislike supera quello dei like.
Uno stesso video può comparire tra i video in tendenza in giornate differenti e dare luogo a più record selezionati: il job restituisce ciascun titolo una sola volta, eliminando i duplicati.
Come richiesto dalla traccia, per il Job 1 vengono riportati risultati parziali: l'elenco completo dei titoli è salvato nel file di output del job, mentre nel report sono riportati il numero di titoli individuati e alcuni esempi.

\subsection{Implementazione}

Il Job 1 è stato implementato attraverso tre classi Java:
\texttt{Job1\_Driver}, \texttt{Job1\_Mapper} e
\texttt{Job1\_Reducer}, utilizzando il framework
Hadoop MapReduce.

\subsubsection{Driver}

La classe \texttt{Job1\_Driver} rappresenta il punto
di ingresso del programma e si occupa della
configurazione e dell'avvio del job MapReduce.

All'interno del metodo \texttt{main()} viene
inizialmente creato un oggetto \texttt{Configuration},
contenente le proprietà di configurazione di Hadoop.

Successivamente, attraverso il metodo
\texttt{Job.getInstance()}, viene istanziato
il job, al quale viene assegnato il nome
\texttt{Job1\_DislikeCounter}.

Il Driver associa al job le classi
\texttt{Job1\_Mapper} e \texttt{Job1\_Reducer}
mediante i metodi \texttt{setMapperClass()}
e \texttt{setReducerClass()}.

Vengono inoltre specificati i tipi delle
coppie chiave-valore prodotte dal Mapper
e dal Reducer, utilizzando le classi
\texttt{Text} e \texttt{NullWritable}.

Il Driver riceve come primo argomento la cartella
del gruppo di dataset e aggiunge come input,
tramite \texttt{FileInputFormat.addInputPath()},
soltanto i file \texttt{*videos.csv} in essa contenuti,
ignorando i JSON utilizzati dal Job 3
(Sezione~\ref{sec:generale}). Se la cartella non contiene
alcun CSV, il programma termina con un messaggio di errore.
Il percorso di output viene impostato con
\texttt{FileOutputFormat.setOutputPath()}.

L'esecuzione del job viene avviata mediante
il metodo \texttt{waitForCompletion(true)},
che attende il completamento delle
operazioni MapReduce; al termine il Driver
stampa il tempo di esecuzione.

\subsubsection{Mapper}

La classe \texttt{Job1\_Mapper} estende la classe
\texttt{Mapper} del framework Hadoop e si occupa
di selezionare, tra i record del dataset, quelli
relativi a video con un numero di dislike
maggiore del numero di like.

Il Mapper riceve in ingresso coppie chiave-valore
di tipo \texttt{LongWritable, Text}.

La chiave, di tipo \texttt{LongWritable},
rappresenta l'offset della riga all'interno del
file di input, mentre il valore, di tipo
\texttt{Text}, contiene il record del dataset.

Per ciascun record, il Mapper esegue le
seguenti operazioni:

\begin{enumerate}
    \item Converte il record ricevuto in una stringa.
    \item Verifica che il record non corrisponda
    all'intestazione del file CSV, riconosciuta dal
    prefisso \texttt{video\_id,}.
    \item Suddivide il record nei campi che lo
    compongono mediante la classe di supporto
    \texttt{CSV\_Parser}, che tiene conto delle virgole
    eventualmente presenti nei campi racchiusi
    tra virgolette.
    \item Verifica che il record contenga almeno dieci
    campi, quelli necessari per accedere alla colonna
    \texttt{dislikes}; in caso contrario il record
    viene scartato.
    \item Estrae il campo \texttt{title} (indice 2),
    il campo \texttt{likes} (indice 8) e il campo
    \texttt{dislikes} (indice 9).
    \item Converte i valori di \texttt{likes} e
    \texttt{dislikes} in numeri interi di tipo
    \texttt{long}, dopo aver eliminato gli spazi iniziali
    e finali; i record per i quali la conversione
    non è possibile vengono ignorati.
    \item Confronta i due valori: se il numero di dislike
    è maggiore del numero di like, emette
    il titolo del video come chiave.
\end{enumerate}

L'output del Mapper è quindi costituito da
coppie di tipo \texttt{Text, NullWritable},
nella forma:

\begin{center}
    \texttt{(TitoloVideo, null)}
\end{center}

Il valore associato alla chiave non è significativo,
poiché l'obiettivo del job consiste nell'ottenere
l'insieme dei titoli e non un valore numerico
a essi associato.

Il Mapper svolge quindi una funzione di filtro:
dei record letti viene emesso soltanto
il sottoinsieme che soddisfa la condizione
\texttt{dislikes > likes}.

\subsubsection{Reducer}

La classe \texttt{Job1\_Reducer} estende la classe
\texttt{Reducer} del framework Hadoop e ha il compito
di produrre l'elenco dei titoli dei video selezionati,
scrivendo ciascun titolo una sola volta.

Il Reducer riceve in ingresso coppie chiave-valore
di tipo \texttt{Text, NullWritable}, raggruppate
automaticamente dal framework Hadoop durante le
fasi di Shuffle e Sort.

La chiave rappresenta il titolo di un video, mentre
la sequenza dei valori contiene un elemento per ciascun
record selezionato dal Mapper con quel titolo.

Un video presente tra i video in tendenza in più
giornate, con dislike superiori ai like, genera
più record e quindi più coppie con la stessa chiave.
Il framework le raggruppa in un unico gruppo,
e il metodo \texttt{reduce()} scrive la chiave
una sola volta, ignorando la molteplicità dei valori.

In questo modo il Reducer elimina i duplicati
senza ricorrere a strutture dati ausiliarie:
la deduplicazione si basa sul raggruppamento
per chiave effettuato dal framework.

Poiché le chiavi vengono ordinate durante la fase
di Sort, il file di output contiene i titoli
in ordine alfabetico, secondo l'ordinamento
dei byte del tipo \texttt{Text}.

Le coppie chiave-valore prodotte dal Reducer
sono di tipo \texttt{Text, NullWritable} e il file
di output contiene un titolo per riga.

\subsection{Pseudocodice}

In questa sezione viene presentato lo pseudocodice
delle tre componenti del Job 1, per descrivere
le principali operazioni eseguite durante
l'elaborazione dei dati.

\subsubsection{Pseudocodice del Driver}

\begin{verbatim}
CLASS Job1_Driver

    FUNCTION main(args)

        IF number of arguments != 2
            PRINT usage message
            EXIT WITH ERROR
        END IF

        startTime = CURRENT_TIME_MILLIS()

        conf = NEW Configuration()

        job = CREATE Job(conf, "Job1_DislikeCounter")

        SET job JAR class = Job1_Driver
        SET job Mapper class = Job1_Mapper
        SET job Reducer class = Job1_Reducer

        SET Mapper output types = (Text, NullWritable)
        SET Reducer output types = (Text, NullWritable)

        FOR EACH file IN directory args[0]
            IF file name ENDS WITH "videos.csv"
                ADD file AS input path
            END IF
        END FOR

        IF no CSV file added
            PRINT error message
            EXIT WITH ERROR
        END IF

        SET output path = args[1]

        success = EXECUTE job AND WAIT

        executionTime = CURRENT_TIME_MILLIS() - startTime

        IF success
            PRINT executionTime
            EXIT SUCCESSFULLY
        ELSE
            PRINT error message
            EXIT WITH ERROR
        END IF

    END FUNCTION

END CLASS
\end{verbatim}

\subsubsection{Pseudocodice del Mapper}

\begin{verbatim}
CLASS Job1_Mapper EXTENDS Mapper

    FUNCTION map(key, value, context)

        line = CONVERT value TO String

        IF line IS CSV HEADER
            RETURN
        END IF

        fields = PARSE CSV line
                 respecting quoted fields

        IF number of fields <= 9
            RETURN
        END IF

        title = fields[2]

        TRY
            likes    = PARSE LONG TRIM(fields[8])
            dislikes = PARSE LONG TRIM(fields[9])
        ON NumberFormatException
            RETURN
        END TRY

        IF dislikes > likes

            EMIT (title, NULL)

        END IF

    END FUNCTION

END CLASS
\end{verbatim}

\subsubsection{Pseudocodice del Reducer}

\begin{verbatim}
CLASS Job1_Reducer EXTENDS Reducer

    FUNCTION reduce(key, values, context)

        // key = video title
        // values = one NULL for each selected record

        EMIT (key, NULL)

    END FUNCTION

END CLASS
\end{verbatim}

\subsection{Pipeline MapReduce}

La pipeline del Job 1 descrive le operazioni
attraverso le quali i record dei dataset YouTube
vengono elaborati per individuare i titoli
dei video con un numero di dislike
maggiore del numero di like.

Il processo si articola nelle seguenti fasi:

\begin{enumerate}

    \item \textbf{Lettura dei dati di input.}

    Hadoop legge i file CSV contenenti i video
    in tendenza e suddivide i dati in porzioni
    denominate \textit{input split}.

    I record vengono forniti ai Mapper sotto forma
    di coppie chiave-valore, dove la chiave
    rappresenta l'offset del record nel file
    e il valore contiene la riga del dataset.

    \item \textbf{Fase di Map.}

    Il Mapper estrae da ciascun record valido
    il titolo, il numero di like e il numero
    di dislike del video.

    Se il numero di dislike è maggiore del numero
    di like, viene emessa una coppia chiave-valore
    nella forma \texttt{(TitoloVideo, null)}.
    In caso contrario il record viene scartato.

    \item \textbf{Fase di Shuffle e Sort.}

    Il framework Hadoop raggruppa automaticamente
    le coppie prodotte dai Mapper in base
    alla chiave, associando a ciascun titolo
    la sequenza dei valori emessi.

    Ad esempio, se lo stesso titolo viene
    selezionato in tre record differenti
    (ad esempio in tre giornate diverse), si ottiene:

    \begin{center}
        \texttt{(TitoloA, [null, null, null])}
    \end{center}

    Le chiavi vengono inoltre ordinate.

    \item \textbf{Fase di Reduce.}

    Per ciascun gruppo, il Reducer scrive
    il titolo una sola volta, indipendentemente
    dal numero di record selezionati
    che lo riguardano.

    \item \textbf{Produzione dell'output.}

    Il risultato viene scritto nel file
    \texttt{part-r-00000}, con un titolo per riga,
    in ordine alfabetico e senza duplicati,
    e successivamente copiato in un file locale
    con \texttt{hadoop fs -getmerge}.

\end{enumerate}

\needspace{13cm}
La Figura~\ref{fig:job1pipeline} mostra le fasi della pipeline
applicate ad alcuni record del file \texttt{USvideos.csv};
nei record letti, $o_i$ indica l'offset della riga nel file.

\begin{figure}[H]
\centering
\pipeline{3.1cm}
{...,BEIRUT | Official Trailer,...,77,17,...\\
...,BEIRUT | Official Trailer,...,326,354,...\\
...,BEIRUT | Official Trailer,...,406,938,...\\
...,MTB Skills - Technical Climb,...,2,3,...\\
...,MTB Skills - Technical Climb,...,2,4,...}
{($o_1$, ...,BEIRUT | Official...,77,17,...)\\
($o_2$, ...,BEIRUT | Official...,326,354,...)\\
($o_3$, ...,BEIRUT | Official...,406,938,...)\\
($o_4$, ...,MTB Skills - Tech...,2,3,...)\\
($o_5$, ...,MTB Skills - Tech...,2,4,...)}
{{\normalfont\itshape $o_1$ scartato: 17 dislike, 77 like}\\
(BEIRUT | Official Trailer, null)\\
(BEIRUT | Official Trailer, null)\\
(MTB Skills - Technical Climb, null)\\
(MTB Skills - Technical Climb, null)}
{(BEIRUT | Official Trailer,\\
\hspace*{1em}[null, null])\\
(MTB Skills - Technical Climb,\\
\hspace*{1em}[null, null])}
{(BEIRUT | Official Trailer, null)\\
(MTB Skills - Technical Climb, null)}
{BEIRUT | Official Trailer\\
MTB Skills - Technical Climb}

\caption{Pipeline del Job 1 su alcuni record di esempio.}
\label{fig:job1pipeline}
\end{figure}


\subsection{Risultati}

Il Job 1 è stato eseguito sui tre gruppi di dataset
descritti nella Sezione~\ref{sec:gruppi}.
Per ciascun gruppo sono stati individuati
i titoli dei video con un numero di dislike
maggiore del numero di like.

I risultati sono rimasti invariati nelle
tre configurazioni Hadoop considerate.

Come richiesto dalla traccia, i risultati
del Job 1 sono presentati in forma parziale:
per ciascun gruppo vengono riportati
i valori dei contatori del job e alcuni
titoli di esempio, mentre l'elenco completo
è contenuto nel file di output.

\subsubsection{Risultati del Dataset A}

Il Dataset A comprende 3 file CSV.
Il job ha letto 139453 righe di input e ha selezionato 1690 record,
pari a circa l'1,21\% del totale.
Questi record corrispondono a 623 titoli distinti.

\begin{center}
\small
\begin{tabular}{|p{0.85\textwidth}|}
\hline
\textbf{Titolo del video (esempi)} \\
\hline
2 Grande Vanilla Lowfat Latte to Go \\
Amy Schumer Goes Behind the Scenes at SNL \\
Amy Schumer Is Ready to Host the Best SNL Yet \\
Anita Latifi - Bleib so (Official Video) \\
Anne Will: Wie antisemitisch ist Deutschland heute ? \\
\hline
\end{tabular}
\end{center}

\subsubsection{Risultati del Dataset B}

Il Dataset B comprende 6 file CSV.
Il job ha letto 278026 righe di input e ha selezionato 4464 record,
pari a circa l'1,61\% del totale.
Questi record corrispondono a 2748 titoli distinti.

\begin{center}
\small
\begin{tabular}{|p{0.85\textwidth}|}
\hline
\textbf{Titolo del video (esempi)} \\
\hline
20 second video \\
A Wardrobe Malfunction Might Cost This Team An Olympic Gold Medal \\
A new Aeroplan \\
ABC cancels 'Roseanne' after Barr's racist tweet \\
Affaire Maëlys : L'avocat de Nordahl Lelandais s'exprime \\
\hline
\end{tabular}
\end{center}

\subsubsection{Risultati del Dataset C}

Il Dataset C comprende 10 file CSV.
Il job ha letto 419221 righe di input e ha selezionato 6886 record,
pari a circa l'1,64\% del totale.
Questi record corrispondono a 4145 titoli distinti.

\begin{center}
\small
\begin{tabular}{|p{0.85\textwidth}|}
\hline
\textbf{Titolo del video (esempi)} \\
\hline
ASI SE BAILA EL VALS EN EL 2018 \\
Accidente en la carretera México - Puebla, deja tres fallecidos \\
Actress Sri Reddy Reacts On Pawan Kalyan Comment \\
Anchor Pradeep's appeal against drunk driving! - TV9 Now \\
Anupama Married! Who is Her Husband? \\
\hline
\end{tabular}
\end{center}

\subsubsection{Confronto dei risultati}

Le grandezze osservate nei tre gruppi di dataset
sono riassunte nella Tabella~\ref{tab:job1risultati}.

\begin{table}[H]
\centering
\small
\begin{tabular}{|l|r|r|r|}
\hline
\textbf{Grandezza} &
\textbf{A} &
\textbf{B} &
\textbf{C} \\
\hline
Righe lette (\textit{Map input}) & 139453 & 278026 & 419221 \\
Record selezionati (\textit{Map output}) & 1690 & 4464 & 6886 \\
Titoli distinti (output del job) & 623 & 2748 & 4145 \\
\hline
\end{tabular}
\caption{Risultati del Job 1 nei tre gruppi di dataset.}
\label{tab:job1risultati}
\end{table}

Il numero di righe lette è superiore al numero
di record effettivi dei file CSV a causa delle
descrizioni distribuite su più righe
(Sezione~\ref{sec:parser}).

Dal confronto emerge che l'elenco dei titoli
individuati nel Dataset A è interamente contenuto
in quello del Dataset B, che a sua volta è contenuto
in quello del Dataset C.
Questo risultato è coerente con la definizione dei gruppi,
ciascuno dei quali comprende il precedente.

L'aggiunta di nuovi Paesi aumenta quindi
il numero di titoli individuati: si passa
da 623 titoli nel Dataset A a 2748 nel Dataset B
e a 4145 nel Dataset C.

La percentuale di record selezionati
rimane contenuta in tutti e tre i gruppi,
tra l'1,21\% e l'1,64\%: i video con più
dislike che like costituiscono quindi
una piccola frazione dei record in tendenza.

In tutti e tre i gruppi il numero di record selezionati
è maggiore del numero di titoli distinti
(ad esempio 1690 contro 623 nel Dataset A).
Questo conferma che numerosi video compaiono
tra i video in tendenza in più giornate
e che la deduplicazione svolta dal Reducer
è necessaria per ottenere un elenco di titoli.

\subsection{Analisi delle prestazioni}

Il Job 1 è stato eseguito nelle tre configurazioni
descritte nella Tabella~\ref{tab:configurazioni},
per ciascuno dei gruppi di dataset A, B e C.

\subsubsection{Esecuzione in modalità Local}

In questa configurazione i dati vengono letti
dal filesystem locale e il job viene eseguito
dal framework MapReduce locale, senza HDFS e YARN.
I tempi ottenuti sono riportati nella
Tabella~\ref{tab:job1local}.

\begin{table}[H]
\centering
\small
\begin{tabular}{|l|r|r|r|}
\hline
\textbf{Esecuzione} &
\textbf{A (ms)} &
\textbf{B (ms)} &
\textbf{C (ms)} \\
\hline
Esecuzione 1 & 10346 & 13551 & 11276 \\
Esecuzione 2 & 5310 & 8627 & 9125 \\
Esecuzione 3 & 5161 & 8330 & 11771 \\
\hline
\textbf{Media} & 6939 & 10169 & 10724 \\
\hline
\end{tabular}
\caption{Tempi di esecuzione del Job 1 in modalità Local.}
\label{tab:job1local}
\end{table}

Il framework ha suddiviso i dati di input
in 6, 13 e 19 \textit{split} rispettivamente
per i dataset A, B e C: sul filesystem locale
i blocchi hanno una dimensione di 32 MB e i file
CSV più grandi vengono quindi divisi in più split.

La prima esecuzione è stata più lenta delle altre due
(ad esempio 10346 millisecondi contro 5310 e 5161
per il Dataset A), per l'avvio a freddo descritto
nella Sezione~\ref{sec:tempi}. Questo effetto alza
le medie del Job 1 in modalità Local, soprattutto
per i Dataset A e B.

\subsubsection{Esecuzione in modalità
Pseudo-distributed senza YARN}

In questa configurazione i dataset sono memorizzati
su HDFS, mentre il job viene eseguito dal framework
MapReduce locale, senza ricorrere al gestore di risorse YARN.
I tempi ottenuti sono riportati nella
Tabella~\ref{tab:job1pd}.

\begin{table}[H]
\centering
\small
\begin{tabular}{|l|r|r|r|}
\hline
\textbf{Esecuzione} &
\textbf{A (ms)} &
\textbf{B (ms)} &
\textbf{C (ms)} \\
\hline
Esecuzione 1 & 7538 & 10338 & 12055 \\
Esecuzione 2 & 6195 & 9225 & 13308 \\
Esecuzione 3 & 6075 & 9170 & 11644 \\
\hline
\textbf{Media} & 6603 & 9578 & 12336 \\
\hline
\end{tabular}
\caption{Tempi di esecuzione del Job 1 in modalità Pseudo-distributed senza YARN.}
\label{tab:job1pd}
\end{table}

Su HDFS i blocchi hanno una dimensione di 128 MB:
poiché tutti i file CSV sono più piccoli, ciascuno
costituisce un unico split e i map task sono stati
3, 6 e 10 rispettivamente per i dataset A, B e C.

\subsubsection{Esecuzione in modalità
Pseudo-distributed con YARN}

Nella terza configurazione HDFS è stato impiegato
per la memorizzazione dei dataset, mentre YARN
ha gestito l'assegnazione delle risorse
e l'esecuzione dei task MapReduce.
I tempi ottenuti sono riportati nella
Tabella~\ref{tab:job1yarn}.

\begin{table}[H]
\centering
\small
\begin{tabular}{|l|r|r|r|}
\hline
\textbf{Esecuzione} &
\textbf{A (ms)} &
\textbf{B (ms)} &
\textbf{C (ms)} \\
\hline
Esecuzione 1 & 54758 & 69533 & 112323 \\
Esecuzione 2 & 36837 & 68131 & 105203 \\
Esecuzione 3 & 35592 & 65044 & 93560 \\
\hline
\textbf{Media} & 42396 & 67569 & 103695 \\
\hline
\end{tabular}
\caption{Tempi di esecuzione del Job 1 in modalità Pseudo-distributed con YARN.}
\label{tab:job1yarn}
\end{table}

Trattandosi di un unico nodo, tutti i map task vengono
eseguiti sul nodo che contiene i dati (\textit{data-local}).
YARN può inoltre avviare una copia dei task più lenti
e conservare il risultato di quella che termina per prima
(\textit{esecuzione speculativa}): le copie scartate compaiono
nei log come \texttt{Killed map tasks}, ma non si tratta
di errori e il risultato del job non cambia.

\subsubsection{Confronto delle configurazioni}

Nella Tabella~\ref{tab:job1confronto} sono riportati
i tempi medi delle tre configurazioni.

\begin{table}[H]
\centering
\small
\begin{tabular}{|l|r|r|r|}
\hline
\textbf{Modalità} &
\textbf{A (ms)} &
\textbf{B (ms)} &
\textbf{C (ms)} \\
\hline
Local & 6939 & 10169 & 10724 \\
PD senza YARN & 6603 & 9578 & 12336 \\
PD con YARN & 42396 & 67569 & 103695 \\
\hline
\end{tabular}
\caption{Tempi medi di esecuzione del Job 1 nelle tre configurazioni.}
\label{tab:job1confronto}
\end{table}

La configurazione con YARN ha richiesto tempi
sensibilmente superiori alle altre due: in media da circa 6,1
a circa 9,7 volte il tempo della modalità Local,
con lo scarto più ampio per il Dataset C.
Su una singola macchina YARN aggiunge il costo
dell'avvio dell'ApplicationMaster e di un container
per ciascun task, senza poter distribuire il lavoro
su più nodi.

Le modalità Local e Pseudo-distributed senza YARN
hanno tempi medi vicini: la seconda è risultata più veloce
di circa il 5\% per i Dataset A e B e più lenta di circa
il 15\% per il Dataset C. Il vantaggio per i Dataset A e B
dipende però dalla prima esecuzione in modalità Local:
considerando soltanto la seconda e la terza, la modalità
Local sarebbe la più veloce per tutti i gruppi (in media
5236, 8478 e 10448 millisecondi, contro 6603, 9578 e 12336).

In tutte le configurazioni il tempo cresce
dal Dataset A al Dataset C.

I file di output prodotti nelle tre configurazioni
sono risultati identici.

\needspace{10cm}
\subsubsection{Rappresentazione grafica delle prestazioni}

La Figura~\ref{fig:job1tempi} riporta i tempi
di esecuzione del Job 1 nelle tre configurazioni
per i dataset A, B e C.

\begin{figure}[H]
\centering

\begin{tikzpicture}
\begin{axis}[
    ybar,
    width=0.95\textwidth,
    height=6cm,
    bar width=12pt,
    ylabel={Tempo medio (ms)},
    xlabel={Gruppo di dataset},
    symbolic x coords={A,B,C},
    xtick=data,
    ymin=0,
    enlarge x limits=0.25,
    legend style={
        at={(0.5,-0.28)},
        anchor=north,
        font=\footnotesize,
        cells={anchor=west},
        legend columns=3,
        column sep=6pt
    },
    grid=major,
    grid style={dashed,gray!30}
]

\addplot coordinates {
    (A,6939)
    (B,10169)
    (C,10724)
};

\addplot coordinates {
    (A,6603)
    (B,9578)
    (C,12336)
};

\addplot coordinates {
    (A,42396)
    (B,67569)
    (C,103695)
};

\legend{
    Local,
    Pseudo-distributed senza YARN,
    Pseudo-distributed con YARN
}

\end{axis}
\end{tikzpicture}

\caption{Confronto dei tempi medi di esecuzione del Job 1 nelle tre configurazioni Hadoop.}
\label{fig:job1tempi}

\end{figure}

Il grafico evidenzia il distacco della configurazione
con YARN e la vicinanza tra le configurazioni Local
e Pseudo-distributed senza YARN.
Tali risultati vanno interpretati nel contesto
dell'ambiente sperimentale, costituito da una singola
macchina virtuale.

\section{Job 2 - Individuazione dei cinque canali YouTube con il maggior numero di video in tendenza}
\subsection{Obiettivo}
Il secondo job ha l'obiettivo di individuare i cinque canali YouTube che presentano il maggior numero di occorrenze all'interno dei dataset relativi ai video in tendenza (\textit{trending videos}).
L'analisi viene effettuata mediante il paradigma Hadoop MapReduce, attraverso tre componenti principali: Mapper, Reducer e Driver.
Il programma considera il campo \texttt{channel\_title} presente nei file CSV e determina il numero di record associati a ciascun canale. Successivamente, seleziona i cinque canali con il maggior numero di occorrenze.
Il conteggio si riferisce ai record del dataset: pertanto, uno stesso video presente tra i video in tendenza in giornate differenti viene conteggiato più volte.

\subsection{Implementazione}

Il Job 2 è stato implementato attraverso tre classi Java:
\texttt{Job2\_Driver}, \texttt{Job2\_Mapper} e
\texttt{Job2\_Reducer}, utilizzando il framework
Hadoop MapReduce.

\subsubsection{Driver}

La classe \texttt{Job2\_Driver} rappresenta il punto
di ingresso del programma e si occupa della
configurazione e dell'avvio del job MapReduce.

All'interno del metodo \texttt{main()} viene
inizialmente creato un oggetto \texttt{Configuration},
contenente le proprietà di configurazione di Hadoop.

Successivamente, attraverso il metodo
\texttt{Job.getInstance()}, viene istanziato
il job, al quale viene assegnato il nome
\texttt{Job2\_TopChannels}.

Il Driver associa al job le classi
\texttt{Job2\_Mapper} e \texttt{Job2\_Reducer}
mediante i metodi \texttt{setMapperClass()}
e \texttt{setReducerClass()}.

Vengono inoltre specificati i tipi delle
coppie chiave-valore prodotte dal Mapper
e dal Reducer, utilizzando le classi
\texttt{Text} e \texttt{IntWritable}.

Come nel Job 1, il Driver aggiunge come input
soltanto i file \texttt{*videos.csv} della cartella
indicata, ignorando i JSON, e imposta il percorso
di output con \texttt{FileOutputFormat.setOutputPath()}.

L'esecuzione del job viene avviata mediante
il metodo \texttt{waitForCompletion(true)},
che attende il completamento delle
operazioni MapReduce; al termine il Driver
stampa il tempo di esecuzione.

\subsubsection{Mapper}

La classe \texttt{Job2\_Mapper} estende la classe
\texttt{Mapper} del framework Hadoop e si occupa
di estrarre, da ogni record del dataset, il nome
del canale YouTube associato al video.

Il Mapper riceve in ingresso coppie chiave-valore
di tipo \texttt{LongWritable, Text}.

La chiave, di tipo \texttt{LongWritable},
rappresenta l'offset della riga all'interno del
file di input, mentre il valore, di tipo
\texttt{Text}, contiene il record del dataset.

Per ciascun record, il Mapper esegue le
seguenti operazioni:

\begin{enumerate}
    \item Converte il record ricevuto in una stringa.
    \item Verifica che il record non corrisponda
    all'intestazione del file CSV.
    \item Suddivide il record nei campi che lo
    compongono, tenendo conto delle virgole
    eventualmente presenti nei campi racchiusi
    tra virgolette.
    \item Estrae il campo \texttt{channel\_title},
    situato nella quarta colonna del dataset
    (indice 3).
    \item Verifica che il nome del canale
    non sia vuoto.
    \item Emette una coppia chiave-valore
    costituita dal nome del canale e dal
    valore intero 1.
\end{enumerate}

L'output del Mapper è quindi costituito da
coppie di tipo \texttt{Text, IntWritable},
nella forma:

\begin{center}
    \texttt{(NomeCanale, 1)}
\end{center}

Ogni coppia rappresenta una singola occorrenza
del canale all'interno del dataset.

Il valore 1 viene assegnato a ogni record valido
indipendentemente dal numero di visualizzazioni
del video, poiché l'obiettivo del job consiste
nel conteggiare le occorrenze dei canali e non
nell'analizzare le visualizzazioni.

\subsubsection{Reducer}

La classe \texttt{Job2\_Reducer} estende la classe
\texttt{Reducer} del framework Hadoop e ha il compito
di determinare i cinque canali YouTube con il maggior
numero di occorrenze nel dataset.

Il Reducer riceve in ingresso coppie chiave-valore
di tipo \texttt{Text, IntWritable}, raggruppate
automaticamente dal framework Hadoop durante le
fasi di Shuffle e Sort.

La chiave rappresenta il nome di un canale YouTube,
mentre i valori costituiscono una sequenza di interi
uguali a 1, ciascuno associato a un record del dataset.

Per ogni canale, il metodo \texttt{reduce()} somma
i valori ricevuti, ottenendo il numero complessivo
di occorrenze del canale.

Per individuare i cinque canali con il maggior
numero di occorrenze viene utilizzata una struttura
dati \texttt{PriorityQueue}, che permette di mantenere
in memoria soltanto i cinque risultati migliori.

La coda viene ordinata in modo da mantenere
in testa il canale con il conteggio minore
tra quelli attualmente selezionati.

Quando viene elaborato un nuovo canale,
il suo conteggio viene confrontato con quello
dell'elemento in testa alla coda.

Se il nuovo conteggio è maggiore, il canale
con il conteggio minore viene rimosso e
sostituito con quello appena elaborato.

In caso di parità di conteggio, viene considerato
peggiore il canale che segue in ordine alfabetico;
lo stesso criterio di spareggio è usato nel metodo
\texttt{cleanup()} per ordinare l'output finale.

Al termine dell'elaborazione, il metodo
\texttt{cleanup()} ordina i cinque canali
selezionati in ordine decrescente di occorrenze
e li scrive nel file di output.

Le coppie chiave-valore prodotte dal Reducer
sono di tipo \texttt{Text, IntWritable}, nella forma:

\begin{center}
    \texttt{(NomeCanale, NumeroOccorrenze)}
\end{center}

Il risultato finale consiste quindi in cinque
coppie chiave-valore, una per ciascuno dei
canali YouTube maggiormente presenti tra
i record dei video in tendenza.

\subsection{Pseudocodice}

In questa sezione viene presentato lo pseudocodice
delle tre componenti del Job 2, per descrivere
le principali operazioni eseguite durante
l'elaborazione dei dati.

\subsubsection{Pseudocodice del Driver}

\begin{verbatim}
CLASS Job2_Driver

    FUNCTION main(args)

        IF number of arguments != 2
            PRINT usage message
            EXIT WITH ERROR
        END IF

        startTime = CURRENT_TIME_MILLIS()

        conf = NEW Configuration()

        job = CREATE Job(conf, "Job2_TopChannels")

        SET job JAR class = Job2_Driver
        SET job Mapper class = Job2_Mapper
        SET job Reducer class = Job2_Reducer

        SET Mapper output types = (Text, IntWritable)
        SET Reducer output types = (Text, IntWritable)

        SET number of Reducers = 1

        FOR EACH file IN directory args[0]
            IF file name ENDS WITH "videos.csv"
                ADD file AS input path
            END IF
        END FOR

        IF no CSV file added
            PRINT error message
            EXIT WITH ERROR
        END IF

        SET output path = args[1]

        success = EXECUTE job AND WAIT

        executionTime = CURRENT_TIME_MILLIS() - startTime

        IF success
            PRINT executionTime
            EXIT SUCCESSFULLY
        ELSE
            PRINT error message
            EXIT WITH ERROR
        END IF

    END FUNCTION

END CLASS
\end{verbatim}

\subsubsection{Pseudocodice del Mapper}

\begin{verbatim}
CLASS Job2_Mapper EXTENDS Mapper

    FUNCTION map(key, value, context)

        line = CONVERT value TO String

        IF line IS CSV HEADER
            RETURN
        END IF

        fields = PARSE CSV line
                 respecting quoted fields

        IF channel_title field EXISTS

            channel = fields[3]

            IF channel IS NOT EMPTY

                EMIT (channel, 1)

            END IF

        END IF

    END FUNCTION

END CLASS
\end{verbatim}

\subsubsection{Pseudocodice del Reducer}

\begin{verbatim}
CLASS Job2_Reducer EXTENDS Reducer

    TOP_K = 5
    INITIALIZE PriorityQueue topChannels
               WITH THE WORST CHANNEL AT THE HEAD

    FUNCTION reduce(key, values, context)

        count = 0

        FOR EACH value IN values
            count = count + value
        END FOR

        IF topChannels.size < TOP_K

            INSERT (key, count) INTO topChannels

        ELSE IF (key, count) IS BETTER THAN
                THE HEAD OF topChannels

            REMOVE HEAD FROM topChannels
            INSERT (key, count) INTO topChannels

        END IF

    END FUNCTION

    FUNCTION cleanup(context)

        ranking = SORT topChannels BY count DESCENDING,
                  THEN BY channel ASCENDING

        FOR EACH element IN ranking

            EMIT (element.channel,
                  element.count)

        END FOR

    END FUNCTION

END CLASS
\end{verbatim}

\subsection{Pipeline MapReduce}

La pipeline del Job 2 descrive le operazioni
attraverso le quali i record dei dataset YouTube
vengono elaborati per individuare i cinque canali
con il maggior numero di occorrenze.

Il processo si articola nelle seguenti fasi:

\begin{enumerate}

    \item \textbf{Lettura dei dati di input.}

    Hadoop legge i file CSV contenenti i video
    in tendenza e suddivide i dati in porzioni
    denominate \textit{input split}.

    I record vengono forniti ai Mapper sotto forma
    di coppie chiave-valore, dove la chiave
    rappresenta l'offset del record nel file
    e il valore contiene la riga del dataset.

    \item \textbf{Fase di Map.}

    Il Mapper estrae da ciascun record valido
    il nome del canale YouTube, contenuto nel
    campo \texttt{channel\_title}.

    Per ogni canale individuato viene emessa
    una coppia chiave-valore nella forma
    \texttt{(NomeCanale, 1)}.

    \item \textbf{Fase di Shuffle e Sort.}

    Il framework Hadoop raggruppa automaticamente
    le coppie prodotte dai Mapper in base
    alla chiave, associando a ciascun canale
    la sequenza dei valori interi emessi.

    Ad esempio, tre occorrenze dello stesso
    canale producono:

    \begin{center}
        \texttt{(CanaleA, [1, 1, 1])}
    \end{center}

    \item \textbf{Fase di Reduce.}

    Per ciascun canale, il Reducer somma
    i valori ricevuti, determinando il numero
    totale di occorrenze nel dataset.

    I conteggi vengono confrontati mediante
    una \texttt{PriorityQueue}, che mantiene
    i cinque canali con il maggior numero
    di occorrenze.

    \item \textbf{Produzione dell'output.}

    Al termine dell'elaborazione, il metodo
    \texttt{cleanup()} ordina i cinque canali
    selezionati in ordine decrescente
    rispetto al numero di occorrenze.

    Il risultato viene scritto nel file
    \texttt{part-r-00000}, sotto forma di
    coppie \texttt{(NomeCanale, NumeroOccorrenze)}.

\end{enumerate}

\needspace{13cm}
La Figura~\ref{fig:job2pipeline} mostra le fasi della pipeline
applicate ad alcuni record del file \texttt{USvideos.csv};
nei record letti, $o_i$ indica l'offset della riga nel file.

\begin{figure}[H]
\centering
\pipeline{3.1cm}
{agKDPntMv-E,17.14.11,...,ESPN,...\\
TUmyygCMMGA,17.14.11,...,Vox,...\\
KJ12FVmumeA,17.15.11,...,ESPN,...\\
zy0b9e40tK8,17.14.11,...,Netflix,...\\
TUmyygCMMGA,17.15.11,...,Vox,...}
{($o_1$, agKDPntMv-E,...,ESPN,...)\\
($o_2$, TUmyygCMMGA,...,Vox,...)\\
($o_3$, KJ12FVmumeA,...,ESPN,...)\\
($o_4$, zy0b9e40tK8,...,Netflix,...)\\
($o_5$, TUmyygCMMGA,...,Vox,...)}
{(ESPN, 1)\\
(Vox, 1)\\
(ESPN, 1)\\
(Netflix, 1)\\
(Vox, 1)}
{(ESPN, [1, 1])\\
(Netflix, [1])\\
(Vox, [1, 1])}
{{\normalfont\itshape reduce():} ESPN=2, Netflix=1, Vox=2\\[3pt]
{\normalfont\itshape cleanup(), primi 5:}\\
(ESPN, 2)\\
(Vox, 2)\\
(Netflix, 1)}
{ESPN\hspace*{2em}2\\
Vox\hspace*{2.6em}2\\
Netflix\hspace*{0.6em}1}

\caption{Pipeline del Job 2 su alcuni record di esempio.}
\label{fig:job2pipeline}
\end{figure}


\subsection{Risultati}

Il Job 2 è stato eseguito sui tre gruppi di dataset
descritti nella Sezione~\ref{sec:gruppi}.
Per ciascun gruppo sono stati individuati
i cinque canali YouTube con il maggior
numero di occorrenze nei record dei video
in tendenza.

I risultati sono rimasti invariati nelle
tre configurazioni Hadoop considerate.

\subsubsection{Risultati del Dataset A}

\begin{center}
\begin{tabular}{|l|r|}
\hline
\textbf{Canale YouTube} & \textbf{Occorrenze} \\
\hline
The Late Show with Stephen Colbert & 502 \\
The Tonight Show Starring Jimmy Fallon & 484 \\
Late Night with Seth Meyers & 475 \\
TheEllenShow & 473 \\
Jimmy Kimmel Live & 465 \\
\hline
\end{tabular}
\end{center}

\subsubsection{Risultati del Dataset B}

\begin{center}
\begin{tabular}{|l|r|}
\hline
\textbf{Canale YouTube} & \textbf{Occorrenze} \\
\hline
The Late Show with Stephen Colbert & 779 \\
Late Night with Seth Meyers & 692 \\
TheEllenShow & 660 \\
The Tonight Show Starring Jimmy Fallon & 653 \\
Jimmy Kimmel Live & 636 \\
\hline
\end{tabular}
\end{center}

\subsubsection{Risultati del Dataset C}

\begin{center}
\begin{tabular}{|l|r|}
\hline
\textbf{Canale YouTube} & \textbf{Occorrenze} \\
\hline
The Late Show with Stephen Colbert & 984 \\
WWE & 804 \\
Late Night with Seth Meyers & 773 \\
VikatanTV & 763 \\
TheEllenShow & 743 \\
\hline
\end{tabular}
\end{center}

\subsubsection{Confronto dei risultati}

Dal confronto emerge che il canale
\textit{The Late Show with Stephen Colbert}
occupa la prima posizione in tutti e tre
i gruppi di dataset.

Passando dal Dataset A al Dataset B,
i cinque canali presenti nella classifica
rimangono gli stessi, ma cambiano
alcune posizioni e aumentano le occorrenze.

Nel Dataset C, invece, entrano nella
classifica i canali \textit{WWE}
e \textit{VikatanTV}, mentre non compaiono
più tra i primi cinque
\textit{The Tonight Show Starring Jimmy Fallon}
e \textit{Jimmy Kimmel Live}.

Questi risultati mostrano come
l'ampliamento dei Paesi considerati
possa modificare la classifica dei canali
maggiormente presenti tra i video
in tendenza.

\subsection{Analisi delle prestazioni}

Il Job 2 è stato eseguito nelle tre configurazioni
descritte nella Tabella~\ref{tab:configurazioni},
per ciascuno dei gruppi di dataset A, B e C.

\subsubsection{Esecuzione in modalità Local}

In questa configurazione i dati vengono letti
dal filesystem locale e il job viene eseguito
dal framework MapReduce locale, senza HDFS e YARN.
I tempi ottenuti sono riportati nella
Tabella~\ref{tab:job2local}.

\begin{table}[H]
\centering
\small
\begin{tabular}{|l|r|r|r|}
\hline
\textbf{Esecuzione} &
\textbf{A (ms)} &
\textbf{B (ms)} &
\textbf{C (ms)} \\
\hline
Esecuzione 1 & 6557 & 9427 & 11221 \\
Esecuzione 2 & 6748 & 8276 & 13554 \\
Esecuzione 3 & 6945 & 9785 & 11392 \\
\hline
\textbf{Media} & 6750 & 9163 & 12056 \\
\hline
\end{tabular}
\caption{Tempi di esecuzione del Job 2 in modalità Local.}
\label{tab:job2local}
\end{table}

Il framework ha suddiviso i dati di input
in 6, 13 e 19 \textit{split} rispettivamente
per i dataset A, B e C: sul filesystem locale
i blocchi hanno una dimensione di 32 MB e i file
CSV più grandi vengono quindi divisi in più split.

\subsubsection{Esecuzione in modalità
Pseudo-distributed senza YARN}

In questa configurazione i dataset sono memorizzati
su HDFS, mentre il job viene eseguito dal framework
MapReduce locale, senza ricorrere al gestore di risorse YARN.
I tempi ottenuti sono riportati nella
Tabella~\ref{tab:job2pd}.

\begin{table}[H]
\centering
\small
\begin{tabular}{|l|r|r|r|}
\hline
\textbf{Esecuzione} &
\textbf{A (ms)} &
\textbf{B (ms)} &
\textbf{C (ms)} \\
\hline
Esecuzione 1 & 7151 & 10240 & 13166 \\
Esecuzione 2 & 7191 & 10158 & 12254 \\
Esecuzione 3 & 7403 & 10475 & 12260 \\
\hline
\textbf{Media} & 7248 & 10291 & 12560 \\
\hline
\end{tabular}
\caption{Tempi di esecuzione del Job 2 in modalità Pseudo-distributed senza YARN.}
\label{tab:job2pd}
\end{table}

Su HDFS i blocchi hanno una dimensione di 128 MB:
poiché tutti i file CSV sono più piccoli, ciascuno
costituisce un unico split e i map task sono stati
3, 6 e 10 rispettivamente per i dataset A, B e C.

\subsubsection{Esecuzione in modalità
Pseudo-distributed con YARN}

Nella terza configurazione HDFS è stato impiegato
per la memorizzazione dei dataset, mentre YARN
ha gestito l'assegnazione delle risorse
e l'esecuzione dei task MapReduce.
I tempi ottenuti sono riportati nella
Tabella~\ref{tab:job2yarn}.

\begin{table}[H]
\centering
\small
\begin{tabular}{|l|r|r|r|}
\hline
\textbf{Esecuzione} &
\textbf{A (ms)} &
\textbf{B (ms)} &
\textbf{C (ms)} \\
\hline
Esecuzione 1 & 44958 & 65952 & 111218 \\
Esecuzione 2 & 45090 & 79543 & 97286 \\
Esecuzione 3 & 49305 & 54103 & 97490 \\
\hline
\textbf{Media} & 46451 & 66533 & 101998 \\
\hline
\end{tabular}
\caption{Tempi di esecuzione del Job 2 in modalità Pseudo-distributed con YARN.}
\label{tab:job2yarn}
\end{table}

Come per il Job 1, tutti i map task sono
\textit{data-local} ed eventuali task interrotti
(\texttt{Killed map tasks}) sono dovuti all'esecuzione
speculativa.

\subsubsection{Confronto delle configurazioni}

Nella Tabella~\ref{tab:job2confronto} sono riportati
i tempi medi delle tre configurazioni.

\begin{table}[H]
\centering
\small
\begin{tabular}{|l|r|r|r|}
\hline
\textbf{Modalità} &
\textbf{A (ms)} &
\textbf{B (ms)} &
\textbf{C (ms)} \\
\hline
Local & 6750 & 9163 & 12056 \\
PD senza YARN & 7248 & 10291 & 12560 \\
PD con YARN & 46451 & 66533 & 101998 \\
\hline
\end{tabular}
\caption{Tempi medi di esecuzione del Job 2 nelle tre configurazioni.}
\label{tab:job2confronto}
\end{table}

Anche per il Job 2 la configurazione con YARN
ha richiesto i tempi più elevati: in media da circa 6,9
a circa 8,5 volte il tempo della modalità Local,
con uno scarto che cresce insieme alla dimensione
dei dati.

La modalità Local è risultata la più veloce in tutti
i gruppi, con tempi medi inferiori a quelli della modalità
Pseudo-distributed senza YARN di circa il 4--11\%.
Quest'ultima è però la più stabile: le sue tre esecuzioni
differiscono al massimo di circa un secondo.
Con tre sole esecuzioni, differenze di questo ordine
restano comparabili con la variabilità delle misure.

A differenza del Job 1, il Mapper del Job 2 emette
una coppia per quasi ogni record: per il Dataset C
sono 379119 coppie \texttt{(NomeCanale, 1)},
tutte trasferite nello shuffle.
Un Combiner che sommi le occorrenze di ciascun canale
all'interno dello stesso map task ridurrebbe
sensibilmente i dati trasferiti nello shuffle.
Il Reducer non può però essere riutilizzato come Combiner,
perché il suo metodo \texttt{cleanup()} seleziona
i cinque canali migliori: servirebbe una classe dedicata
che si limiti alla somma.

I file di output prodotti nelle tre configurazioni
sono risultati identici.

\needspace{10cm}
\subsubsection{Rappresentazione grafica delle prestazioni}

La Figura~\ref{fig:job2tempi} riporta i tempi
di esecuzione del Job 2 nelle tre configurazioni
per i dataset A, B e C.

\begin{figure}[H]
\centering

\begin{tikzpicture}
\begin{axis}[
    ybar,
    width=0.95\textwidth,
    height=6cm,
    bar width=12pt,
    ylabel={Tempo medio (ms)},
    xlabel={Gruppo di dataset},
    symbolic x coords={A,B,C},
    xtick=data,
    ymin=0,
    enlarge x limits=0.25,
    legend style={
        at={(0.5,-0.28)},
        anchor=north,
        font=\footnotesize,
        cells={anchor=west},
        legend columns=3,
        column sep=6pt
    },
    grid=major,
    grid style={dashed,gray!30}
]

\addplot coordinates {
    (A,6750)
    (B,9163)
    (C,12056)
};

\addplot coordinates {
    (A,7248)
    (B,10291)
    (C,12560)
};

\addplot coordinates {
    (A,46451)
    (B,66533)
    (C,101998)
};

\legend{
    Local,
    Pseudo-distributed senza YARN,
    Pseudo-distributed con YARN
}

\end{axis}
\end{tikzpicture}

\caption{Confronto dei tempi medi di esecuzione del Job 2 nelle tre configurazioni Hadoop.}
\label{fig:job2tempi}

\end{figure}

Il grafico evidenzia il distacco della configurazione
con YARN, che cresce dal Dataset A al Dataset C,
e la vicinanza tra le configurazioni Local
e Pseudo-distributed senza YARN.

\section{Job 3 - Individuazione della categoria con il maggior numero medio di visualizzazioni}
\label{sec:job3}

\subsection{Obiettivo}

Il terzo job ha l'obiettivo di individuare il titolo
della categoria di video con il maggior numero medio
di visualizzazioni (\textit{views}).

Nei file CSV ogni video è associato soltanto
a un identificativo numerico, \texttt{category\_id};
il titolo della categoria si trova invece nei file JSON,
uno per ciascun Paese. Il job deve quindi unire
(\textit{join}) le informazioni dei CSV con quelle dei JSON.

La media è calcolata sui record del dataset:
come nel Job 2, un video presente tra i video in tendenza
in più giornate contribuisce con ciascuno dei suoi record,
ognuno con il numero di visualizzazioni registrato quel giorno.

\subsection{Scelte progettuali}

Il Job 3 è il più articolato dei tre, perché richiede
dati provenienti da due tipi di file. Prima
dell'implementazione sono stati quindi analizzati
i file JSON e valutate diverse soluzioni.

\subsubsection{Analisi dei file JSON}

Ogni file JSON contiene l'elenco delle categorie
di un Paese, nella forma (estratto):

\begin{verbatim}
{
 "kind": "youtube#videoCategoryListResponse",
 "items": [
  {
   "kind": "youtube#videoCategory",
   "id": "10",
   "snippet": {
    "channelId": "UCBR8-60-B28hp2BmDPdntcQ",
    "title": "Music",
    "assignable": true
   }
  },
  ...
\end{verbatim}

Per il job sono rilevanti soltanto i campi \texttt{id}
e \texttt{snippet.title}. L'analisi dei dieci file ha evidenziato che:

\begin{itemize}
    \item uno stesso identificativo corrisponde sempre
    allo stesso titolo in tutti i Paesi;
    \item dei 32 identificativi presenti nei JSON,
    nei CSV ne vengono utilizzati soltanto 18;
    \item l'identificativo 29 (\textit{Nonprofits \& Activism})
    compare soltanto nel JSON degli Stati Uniti, ma è utilizzato
    nei CSV di tutti i Paesi: negli altri nove ci sono
    2638 record con questa categoria, di cui 1541 in Russia.
\end{itemize}

Per questo motivo il job non utilizza il JSON del singolo
Paese, ma costruisce un unico dizionario a partire da tutti
i JSON presenti nella cartella del gruppo.
Poiché gli Stati Uniti fanno parte di tutti e tre i gruppi,
anche l'identificativo 29 viene sempre tradotto; inoltre
il Mapper non ha bisogno di sapere da quale Paese
proviene ciascun record.

Utilizzando invece il JSON di ciascun Paese, i 2638 record
sarebbero stati scartati: la categoria vincente non
sarebbe cambiata, ma \textit{Nonprofits \& Activism},
rappresentata dai soli 57 record statunitensi,
sarebbe salita al secondo posto in tutti i gruppi.

\subsubsection{Il join tra CSV e JSON: la Distributed Cache}

I file JSON sono molto piccoli (circa 8 KB ciascuno),
mentre i CSV occupano centinaia di megabyte.
Si è quindi scelto un \textit{map-side join}: il Driver
registra i JSON nella \textit{Distributed Cache} di Hadoop
con il metodo \texttt{job.addCacheFile()}, ogni map task
li legge per intero nel metodo \texttt{setup()} e li carica
in una \texttt{HashMap}, e il metodo \texttt{map()}
traduce l'identificativo di ciascun record nel titolo
della categoria. In questo modo il join non richiede
un job aggiuntivo e i dati dei JSON non passano dallo shuffle.

Le alternative considerate sono riassunte
nella Tabella~\ref{tab:job3alternative}.

\begin{table}[H]
\centering
\small
\begin{tabular}{|p{0.32\textwidth}|p{0.58\textwidth}|}
\hline
\textbf{Alternativa} & \textbf{Motivo per cui è stata scartata} \\
\hline
\textit{Reduce-side join}: CSV e JSON entrambi in input,
unione nel Reducer &
Tutti i record passerebbero dallo shuffle e servirebbe
un secondo job per raggruppare per titolo; inoltre i JSON,
distribuiti su più righe, non si prestano a essere
letti riga per riga come input. \\
\hline
Dizionario passato tramite la \texttt{Configuration} &
Adatto soltanto a dati minuscoli e sposta nel Driver
la lettura dei JSON. \\
\hline
Dizionario scritto nel codice &
Non utilizza i file JSON, come invece richiesto
dalla traccia. \\
\hline
\end{tabular}
\caption{Alternative valutate per il join tra CSV e JSON.}
\label{tab:job3alternative}
\end{table}

Nel metodo \texttt{setup()} i file vengono aperti tramite
il \texttt{FileSystem} associato al loro indirizzo:
la stessa istruzione funziona quindi sia sul filesystem
locale sia su HDFS.

\subsubsection{Lettura dei file JSON}

Per leggere i JSON non è stata utilizzata una libreria
dedicata, che avrebbe aggiunto una dipendenza al jar.
Nei file del dataset ogni chiave occupa una riga distinta
e il campo \texttt{"title"} segue sempre il proprio
\texttt{"id"}: è quindi sufficiente leggere il file riga
per riga, ricordare l'ultimo identificativo incontrato
e associarlo al titolo successivo.
La soluzione è semplice ma dipende dalla formattazione
dei file: un JSON scritto su una sola riga richiederebbe
un vero parser.

\subsubsection{Calcolo della media e tipi numerici}

La media non è un'operazione associativa: la media
di medie parziali è in generale diversa dalla media
complessiva. Il Reducer calcola quindi, per ciascuna
categoria, la somma delle visualizzazioni e il numero
di record, ed effettua la divisione soltanto alla fine.

I valori in gioco sono elevati: il singolo record più visto
supera i 424 milioni di visualizzazioni e la somma
complessiva per il Dataset C è di circa $5 \cdot 10^{11}$,
ben oltre il limite di un \texttt{int} (circa $2{,}1 \cdot 10^{9}$).
Per questo le visualizzazioni e la loro somma sono gestite
come \texttt{long}. La media è calcolata come \texttt{double},
che rappresenta esattamente gli interi fino a circa
$9 \cdot 10^{15}$, e viene troncata a due cifre decimali
per rendere l'output leggibile.

\subsubsection{Aggregazione locale}

Il Mapper emette una coppia per ciascun record valido:
per il Dataset C sono 375941 coppie, tutte trasferite
nello shuffle, pur essendoci soltanto 18 chiavi distinte.
Aggregando la somma e il conteggio di ogni categoria
all'interno di ciascun map task, ogni task emetterebbe
al massimo 18 coppie.
Il Reducer non può essere riutilizzato come Combiner,
perché restituisce una media e scrive il vincitore nel metodo
\texttt{cleanup()}: servirebbero una classe dedicata e un tipo
\texttt{Writable} per la coppia (somma, conteggio).
Per coerenza con gli altri job è stato mantenuto lo schema
classico (Sezione~\ref{sec:generale}).

\subsection{Implementazione}

Il Job 3 è stato implementato attraverso tre classi Java:
\texttt{Job3\_Driver}, \texttt{Job3\_Mapper} e
\texttt{Job3\_Reducer}, oltre alla classe di supporto
\texttt{CSV\_Parser} condivisa con gli altri job.

\subsubsection{Driver}

La classe \texttt{Job3\_Driver} rappresenta il punto
di ingresso del programma: crea il job, al quale
assegna il nome \texttt{Job3\_MaxAverageViews}, e gli associa
le classi \texttt{Job3\_Mapper} e \texttt{Job3\_Reducer}
e i tipi delle coppie prodotte (\texttt{Text, LongWritable}
per il Mapper e \texttt{Text, DoubleWritable} per il Reducer).

Il Driver scorre la cartella del gruppo di dataset:
i file \texttt{*videos.csv} vengono aggiunti come input
del job, mentre i file \texttt{*\_category\_id.json}
vengono registrati nella Distributed Cache.
Se nella cartella mancano i CSV o i JSON, il programma
termina con un messaggio di errore.
Come negli altri job, l'esecuzione viene avviata con
\texttt{waitForCompletion(true)} e al termine
viene stampato il tempo di esecuzione.

\subsubsection{Mapper}

La classe \texttt{Job3\_Mapper} estende la classe
\texttt{Mapper} del framework Hadoop e si occupa
di associare a ogni record il titolo della sua categoria
e il relativo numero di visualizzazioni.
Il Mapper riceve in ingresso coppie chiave-valore
di tipo \texttt{LongWritable, Text}, con l'offset della riga
come chiave e la riga del dataset come valore.

Il metodo \texttt{setup()}, eseguito una sola volta
per ciascun map task prima della lettura dei record,
legge i file JSON ricevuti tramite la Distributed Cache
e costruisce il dizionario \texttt{category\_id} $\rightarrow$
titolo della categoria.

Per ciascun record, il metodo \texttt{map()} esegue
le seguenti operazioni:

\begin{enumerate}
    \item Verifica che il record non corrisponda
    all'intestazione del file CSV.
    \item Suddivide il record nei campi mediante
    la classe \texttt{CSV\_Parser}.
    \item Verifica che il record contenga almeno otto campi,
    quelli necessari per accedere alla colonna \texttt{views};
    in caso contrario il record viene scartato.
    \item Converte il campo \texttt{category\_id} (indice 4)
    in un \texttt{int} e il campo \texttt{views} (indice 7)
    in un \texttt{long}; i record per i quali la conversione
    non è possibile, o che contengono valori negativi,
    vengono ignorati.
    \item Traduce l'identificativo nel titolo della categoria
    tramite il dizionario; se l'identificativo non è presente
    in alcun JSON, il record viene scartato.
    \item Emette il titolo della categoria come chiave
    e il numero di visualizzazioni come valore.
\end{enumerate}

L'output del Mapper è quindi costituito da coppie
di tipo \texttt{Text, LongWritable}, nella forma:

\begin{center}
    \texttt{(TitoloCategoria, Visualizzazioni)}
\end{center}

\subsubsection{Reducer}

La classe \texttt{Job3\_Reducer} estende la classe
\texttt{Reducer} e ha il compito di individuare
la categoria con la media di visualizzazioni più alta.

Il Reducer riceve, per ciascuna categoria, la sequenza
delle visualizzazioni di tutti i suoi record, raggruppate
dal framework durante le fasi di Shuffle e Sort.
Per ogni categoria il metodo \texttt{reduce()} calcola
la somma delle visualizzazioni e il numero di record,
ne ricava la media e la confronta con la migliore
trovata fino a quel momento, conservando la maggiore.
Poiché le chiavi arrivano in ordine alfabetico e il
confronto è stretto, a parità di media prevale
la categoria che viene prima in ordine alfabetico.

Il metodo \texttt{reduce()} non scrive nulla in output:
è il metodo \texttt{cleanup()}, eseguito una sola volta
al termine dell'elaborazione, a scrivere la categoria
vincente con la sua media, troncata a due cifre decimali.
Lo schema è lo stesso del Job 2, con una classifica
di un solo elemento.
Le coppie prodotte sono di tipo \texttt{Text, DoubleWritable},
nella forma:

\begin{center}
    \texttt{(TitoloCategoria, MediaVisualizzazioni)}
\end{center}

\subsection{Pseudocodice}

\subsubsection{Pseudocodice del Driver}

\begin{verbatim}
CLASS Job3_Driver

    FUNCTION main(args)

        IF number of arguments != 2
            PRINT usage message
            EXIT WITH ERROR
        END IF

        startTime = CURRENT_TIME_MILLIS()

        conf = NEW Configuration()

        job = CREATE Job(conf, "Job3_MaxAverageViews")

        SET job JAR class = Job3_Driver
        SET job Mapper class = Job3_Mapper
        SET job Reducer class = Job3_Reducer

        SET Mapper output types = (Text, LongWritable)
        SET Reducer output types = (Text, DoubleWritable)

        SET number of Reducers = 1

        FOR EACH file IN directory args[0]
            IF file name ENDS WITH "videos.csv"
                ADD file AS input path
            ELSE IF file name ENDS WITH "_category_id.json"
                ADD file TO DISTRIBUTED CACHE
            END IF
        END FOR

        IF no CSV file OR no JSON file added
            PRINT error message
            EXIT WITH ERROR
        END IF

        SET output path = args[1]

        success = EXECUTE job AND WAIT

        executionTime = CURRENT_TIME_MILLIS() - startTime

        IF success
            PRINT executionTime
            EXIT SUCCESSFULLY
        ELSE
            PRINT error message
            EXIT WITH ERROR
        END IF

    END FUNCTION

END CLASS
\end{verbatim}

\subsubsection{Pseudocodice del Mapper}

\begin{verbatim}
CLASS Job3_Mapper EXTENDS Mapper

    categoryNames = EMPTY MAP   // category_id -> title

    FUNCTION setup(context)

        FOR EACH file IN DISTRIBUTED CACHE
            currentId = NULL
            FOR EACH line IN file
                IF line STARTS WITH "id"
                    currentId = VALUE OF line
                ELSE IF line STARTS WITH "title"
                        AND currentId != NULL
                    categoryNames[currentId] = VALUE OF line
                    currentId = NULL
                END IF
            END FOR
        END FOR

    END FUNCTION

    FUNCTION map(key, value, context)

        line = CONVERT value TO String

        IF line IS CSV HEADER
            RETURN
        END IF

        fields = PARSE CSV line
                 respecting quoted fields

        IF number of fields <= 7
            RETURN
        END IF

        TRY
            categoryId = PARSE INT  TRIM(fields[4])
            views      = PARSE LONG TRIM(fields[7])
        ON NumberFormatException
            RETURN
        END TRY

        IF categoryId < 0 OR views < 0
            RETURN
        END IF

        name = categoryNames[categoryId]

        IF name IS NULL
            RETURN
        END IF

        EMIT (name, views)

    END FUNCTION

END CLASS
\end{verbatim}

\subsubsection{Pseudocodice del Reducer}

\begin{verbatim}
CLASS Job3_Reducer EXTENDS Reducer

    bestCategory = NULL
    bestAverage  = -1

    FUNCTION reduce(key, values, context)

        sum = 0
        count = 0

        FOR EACH value IN values
            sum = sum + value
            count = count + 1
        END FOR

        average = sum / count

        IF average > bestAverage
            bestAverage  = average
            bestCategory = key
        END IF

    END FUNCTION

    FUNCTION cleanup(context)

        IF bestCategory != NULL
            EMIT (bestCategory,
                  TRUNCATE bestAverage TO 2 DECIMALS)
        END IF

    END FUNCTION

END CLASS
\end{verbatim}

\subsection{Pipeline MapReduce}

La pipeline del Job 3 si articola nelle seguenti fasi:

\begin{enumerate}

    \item \textbf{Distribuzione dei file JSON.}

    Il Driver registra i file JSON nella Distributed Cache:
    Hadoop li rende disponibili a tutti i map task,
    su qualunque nodo vengano eseguiti.

    \item \textbf{Lettura dei dati di input.}

    Hadoop legge i file CSV e li suddivide in
    \textit{input split}; i record vengono forniti
    ai Mapper come coppie (offset, riga).

    \item \textbf{Preparazione del Mapper.}

    Ogni map task, nel metodo \texttt{setup()},
    costruisce il dizionario delle categorie
    a partire dai file JSON.

    \item \textbf{Fase di Map.}

    Per ciascun record valido viene emessa
    una coppia \texttt{(TitoloCategoria, Visualizzazioni)},
    ad esempio \texttt{(Music, 687582)}.

    \item \textbf{Fase di Shuffle e Sort.}

    Il framework raggruppa le coppie per categoria
    e ordina le chiavi:

    \begin{center}
        \texttt{(Music, [687582, 17158531, 13917, ...])}
    \end{center}

    \item \textbf{Fase di Reduce.}

    Per ciascuna categoria il Reducer calcola la media
    delle visualizzazioni e conserva la maggiore.

    \item \textbf{Produzione dell'output.}

    Il metodo \texttt{cleanup()} scrive un'unica riga
    con la categoria vincente e la sua media nel file
    \texttt{part-r-00000}, successivamente copiato
    in un file locale con \texttt{hadoop fs -getmerge}.

\end{enumerate}

\needspace{13cm}
La Figura~\ref{fig:job3pipeline} mostra le fasi della pipeline
applicate ad alcuni record del file \texttt{USvideos.csv};
nei record letti, $o_i$ indica l'offset della riga nel file.

\begin{figure}[H]
\centering
\pipeline{3.4cm}
{1ZAPwfrtAFY,...,24,...,2418783,...\\
puqaWrEC7tY,...,24,...,343168,...\\
d380meD0W0M,...,24,...,2095731,...\\
5E4ZBSInqUU,...,10,...,687582,...\\
n1WpP7iowLc,...,10,...,17158531,...\\
e\_7zHm7GsYc,...,10,...,13917,...}
{($o_1$, 1ZAPwfrtAFY,...,24,...,2418783,...)\\
($o_2$, puqaWrEC7tY,...,24,...,343168,...)\\
($o_3$, d380meD0W0M,...,24,...,2095731,...)\\
($o_4$, 5E4ZBSInqUU,...,10,...,687582,...)\\
($o_5$, n1WpP7iowLc,...,10,...,17158531,...)\\
($o_6$, e\_7zHm7GsYc,...,10,...,13917,...)}
{{\normalfont\itshape setup(): 10 = Music, 24 = Entertainment}\\
(Entertainment, 2418783)\\
(Entertainment, 343168)\\
(Entertainment, 2095731)\\
(Music, 687582)\\
(Music, 17158531)\\
(Music, 13917)}
{(Entertainment,\\
\hspace*{1em}[2418783, 343168, 2095731])\\
(Music,\\
\hspace*{1em}[687582, 17158531, 13917])}
{{\normalfont\itshape reduce():}\\
Entertainment: 4857682 / 3 = 1619227.33\\
Music: 17860030 / 3 = 5953343.33\\[3pt]
{\normalfont\itshape cleanup(), media massima:}\\
(Music, 5953343.33)}
{Music\hspace*{1em}5953343.33}

\caption{Pipeline del Job 3 su alcuni record di esempio.}
\label{fig:job3pipeline}
\end{figure}


\subsection{Risultati}

Il Job 3 è stato eseguito sui tre gruppi di dataset
descritti nella Sezione~\ref{sec:gruppi}.
I risultati sono rimasti invariati nelle tre
configurazioni Hadoop considerate e sono riassunti
nella Tabella~\ref{tab:job3risultati}.

\begin{table}[H]
\centering
\small
\begin{tabular}{|l|r|r|r|}
\hline
\textbf{Grandezza} &
\textbf{A} &
\textbf{B} &
\textbf{C} \\
\hline
Righe lette (\textit{Map input}) & 139453 & 278026 & 419221 \\
Record validi (\textit{Map output}) & 120705 & 243049 & 375941 \\
Categorie distinte (\textit{Reduce input groups}) & 18 & 18 & 18 \\
\hline
Categoria con la media più alta & Music & Music & Music \\
Media delle visualizzazioni & 9601629,68 & 7360559,37 & 6020771,72 \\
\hline
\end{tabular}
\caption{Risultati del Job 3 nei tre gruppi di dataset.}
\label{tab:job3risultati}
\end{table}

La categoria \textit{Music} ha la media di visualizzazioni
più alta in tutti e tre i gruppi.
La media si riduce passando dal Dataset A al Dataset C:
poiché ogni gruppo comprende il precedente, i video musicali
dei Paesi aggiunti hanno in media meno visualizzazioni
di quelli di Stati Uniti, Regno Unito e Germania.

In tutti i gruppi sono presenti 18 categorie, compresa
\textit{Nonprofits \& Activism}, grazie al dizionario
costruito a partire da tutti i JSON.


\subsection{Analisi delle prestazioni}

Il Job 3 è stato eseguito nelle tre configurazioni
descritte nella Tabella~\ref{tab:configurazioni},
per ciascuno dei gruppi di dataset A, B e C.

\subsubsection{Esecuzione in modalità Local}

In questa configurazione i dati vengono letti
dal filesystem locale e il job viene eseguito
dal framework MapReduce locale, senza HDFS e YARN.
I tempi ottenuti sono riportati nella
Tabella~\ref{tab:job3local}.

\begin{table}[H]
\centering
\small
\begin{tabular}{|l|r|r|r|}
\hline
\textbf{Esecuzione} &
\textbf{A (ms)} &
\textbf{B (ms)} &
\textbf{C (ms)} \\
\hline
Esecuzione 1 & 6798 & 10138 & 13623 \\
Esecuzione 2 & 7341 & 11872 & 12297 \\
Esecuzione 3 & 6797 & 10299 & 16203 \\
\hline
\textbf{Media} & 6979 & 10770 & 14041 \\
\hline
\end{tabular}
\caption{Tempi di esecuzione del Job 3 in modalità Local.}
\label{tab:job3local}
\end{table}

Come per gli altri job, i dati di input sono stati
suddivisi in 6, 13 e 19 \textit{split}
rispettivamente per i dataset A, B e C.

\subsubsection{Esecuzione in modalità
Pseudo-distributed senza YARN}

In questa configurazione i dataset sono memorizzati
su HDFS, mentre il job viene eseguito dal framework
MapReduce locale, senza ricorrere al gestore di risorse YARN.
I tempi ottenuti sono riportati nella
Tabella~\ref{tab:job3pd}.

\begin{table}[H]
\centering
\small
\begin{tabular}{|l|r|r|r|}
\hline
\textbf{Esecuzione} &
\textbf{A (ms)} &
\textbf{B (ms)} &
\textbf{C (ms)} \\
\hline
Esecuzione 1 & 7685 & 11240 & 13253 \\
Esecuzione 2 & 7745 & 10895 & 13394 \\
Esecuzione 3 & 8729 & 12313 & 15853 \\
\hline
\textbf{Media} & 8053 & 11483 & 14167 \\
\hline
\end{tabular}
\caption{Tempi di esecuzione del Job 3 in modalità Pseudo-distributed senza YARN.}
\label{tab:job3pd}
\end{table}

I map task sono stati 3, 6 e 10, uno per ciascun file CSV.
Ogni map task legge inoltre i file JSON nel metodo
\texttt{setup()}: trattandosi di pochi kilobyte,
il costo di questa operazione è trascurabile.

\subsubsection{Esecuzione in modalità
Pseudo-distributed con YARN}

Nella terza configurazione HDFS è stato impiegato
per la memorizzazione dei dataset, mentre YARN
ha gestito l'assegnazione delle risorse
e l'esecuzione dei task MapReduce.
I tempi ottenuti sono riportati nella
Tabella~\ref{tab:job3yarn}.

\begin{table}[H]
\centering
\small
\begin{tabular}{|l|r|r|r|}
\hline
\textbf{Esecuzione} &
\textbf{A (ms)} &
\textbf{B (ms)} &
\textbf{C (ms)} \\
\hline
Esecuzione 1 & 39191 & 53287 & 115340 \\
Esecuzione 2 & 31782 & 78556 & 79873 \\
Esecuzione 3 & 44513 & 49939 & 74831 \\
\hline
\textbf{Media} & 38495 & 60594 & 90015 \\
\hline
\end{tabular}
\caption{Tempi di esecuzione del Job 3 in modalità Pseudo-distributed con YARN.}
\label{tab:job3yarn}
\end{table}

Come per il Job 1, tutti i map task sono
\textit{data-local} ed eventuali task interrotti
(\texttt{Killed map tasks}) sono dovuti all'esecuzione
speculativa.

\subsubsection{Confronto delle configurazioni}

Nella Tabella~\ref{tab:job3confronto} sono riportati
i tempi medi delle tre configurazioni.

\begin{table}[H]
\centering
\small
\begin{tabular}{|l|r|r|r|}
\hline
\textbf{Modalità} &
\textbf{A (ms)} &
\textbf{B (ms)} &
\textbf{C (ms)} \\
\hline
Local & 6979 & 10770 & 14041 \\
PD senza YARN & 8053 & 11483 & 14167 \\
PD con YARN & 38495 & 60594 & 90015 \\
\hline
\end{tabular}
\caption{Tempi medi di esecuzione del Job 3 nelle tre configurazioni.}
\label{tab:job3confronto}
\end{table}

Anche per il Job 3 la configurazione con YARN
ha richiesto i tempi più elevati: in media da circa 5,5
a circa 6,4 volte il tempo della modalità Local, con uno
scarto che cresce insieme alla dimensione dei dati.
È anche la configurazione con la maggiore variabilità:
per il Dataset C le tre esecuzioni vanno da 74831
a 115340 millisecondi, la più lunga dell'intero progetto.

Le modalità Local e Pseudo-distributed senza YARN
hanno tempi medi vicini: la seconda è risultata più lenta
di circa il 15\% per il Dataset A e il 7\% per il Dataset B,
mentre per il Dataset C i tempi coincidono (differenza
inferiore all'1\%).
In tutte le configurazioni il tempo cresce
dal Dataset A al Dataset C.

I file di output prodotti nelle tre configurazioni
sono risultati identici.

\needspace{10cm}
\subsubsection{Rappresentazione grafica delle prestazioni}

La Figura~\ref{fig:job3tempi} riporta i tempi
di esecuzione del Job 3 nelle tre configurazioni
per i dataset A, B e C.

\begin{figure}[H]
\centering

\begin{tikzpicture}
\begin{axis}[
    ybar,
    width=0.95\textwidth,
    height=6cm,
    bar width=12pt,
    ylabel={Tempo medio (ms)},
    xlabel={Gruppo di dataset},
    symbolic x coords={A,B,C},
    xtick=data,
    ymin=0,
    enlarge x limits=0.25,
    legend style={
        at={(0.5,-0.28)},
        anchor=north,
        font=\footnotesize,
        cells={anchor=west},
        legend columns=3,
        column sep=6pt
    },
    grid=major,
    grid style={dashed,gray!30}
]

\addplot coordinates {
    (A,6979)
    (B,10770)
    (C,14041)
};

\addplot coordinates {
    (A,8053)
    (B,11483)
    (C,14167)
};

\addplot coordinates {
    (A,38495)
    (B,60594)
    (C,90015)
};

\legend{
    Local,
    Pseudo-distributed senza YARN,
    Pseudo-distributed con YARN
}

\end{axis}
\end{tikzpicture}

\caption{Confronto dei tempi medi di esecuzione del Job 3 nelle tre configurazioni Hadoop.}
\label{fig:job3tempi}

\end{figure}

Il grafico conferma l'andamento osservato
per i Job 1 e 2: le configurazioni Local
e Pseudo-distributed senza YARN hanno tempi vicini,
mentre quella con YARN se ne discosta sempre di più
al crescere dei dati.

\section{Confronto tra i job e considerazioni finali}

Nella Tabella~\ref{tab:confrontojob} sono riportati
i tempi medi di tutte le combinazioni, calcolati
sulle 81 esecuzioni complessive.

\begin{table}[H]
\centering
\small
\begin{tabular}{|l|l|r|r|r|}
\hline
\textbf{Modalità} & \textbf{Job} &
\textbf{A (ms)} & \textbf{B (ms)} & \textbf{C (ms)} \\
\hline
Local & Job 1 & 6939 & 10169 & 10724 \\
 & Job 2 & 6750 & 9163 & 12056 \\
 & Job 3 & 6979 & 10770 & 14041 \\
\hline
PD senza YARN & Job 1 & 6603 & 9578 & 12336 \\
 & Job 2 & 7248 & 10291 & 12560 \\
 & Job 3 & 8053 & 11483 & 14167 \\
\hline
PD con YARN & Job 1 & 42396 & 67569 & 103695 \\
 & Job 2 & 46451 & 66533 & 101998 \\
 & Job 3 & 38495 & 60594 & 90015 \\
\hline
\end{tabular}
\caption{Tempi medi di esecuzione dei tre job nelle tre configurazioni.}
\label{tab:confrontojob}
\end{table}

\paragraph{Peso dello shuffle.}
Nelle configurazioni Local e Pseudo-distributed senza YARN
i tre job hanno tempi dello stesso ordine sullo stesso gruppo
di dataset: per il Dataset C in modalità Local, ad esempio,
si va da 10724 a 14041 millisecondi.
Eppure i Job 2 e 3 emettono circa 55 volte più coppie
del Job 1 (379119 e 375941 contro 6886).
Con queste dimensioni dei dati, quindi, il tempo
è dominato dalla lettura e dall'analisi dei file CSV
più che dal volume dei dati trasferiti nello shuffle.

\paragraph{Local e Pseudo-distributed senza YARN.}
Le due configurazioni che non usano YARN hanno tempi
molto vicini: la modalità Local è risultata in media
leggermente più veloce in sette casi su nove, con differenze
fino a circa il 15\%, mentre la modalità Pseudo-distributed
senza YARN è stata più veloce soltanto per il Job 1
sui Dataset A e B. Con queste dimensioni dei dati,
quindi, la lettura da HDFS introduce al più
un costo contenuto.

\paragraph{YARN.}
Con YARN i tempi medi per il Dataset C superano
i 90 secondi per tutti e tre i job (da 90015
a 103695 millisecondi). Le differenze tra i job sono
però dello stesso ordine della variabilità tra le singole
esecuzioni, che con YARN arriva fino a circa il 25\%
della media, e non permettono di stabilire quale job
sia il più costoso.
In tutti i casi il costo di avvio dell'ApplicationMaster
e dei container rende YARN la configurazione più lenta
su una singola macchina; i suoi vantaggi emergerebbero
soltanto distribuendo il lavoro su più nodi.

\paragraph{Contatori da interpretare con cautela.}
In modalità Pseudo-distributed senza YARN il contatore
\texttt{HDFS: Number of bytes read} riporta valori superiori
alla dimensione dei file letti (circa 547 MB per i 179 MB
di CSV del Dataset A). Il framework locale esegue infatti
tutti i task nella stessa JVM e le statistiche del filesystem
si accumulano da un task all'altro: il contatore non è quindi
stato utilizzato come misura dei dati letti.

\appendix

\section{Codice sorgente}

Si riporta il codice sorgente dei tre job, contenuto
nel package \texttt{com.BigData} del progetto Maven
\texttt{Esonero\_Big\_Data}.

\subsection{CSV\_Parser}
\lstinputlisting{../Esonero_Big_Data/src/main/java/com/BigData/CSV_Parser.java}

\subsection{Job 1}
\lstinputlisting[title=Job1\_Driver.java]{../Esonero_Big_Data/src/main/java/com/BigData/Job1_Driver.java}
\lstinputlisting[title=Job1\_Mapper.java]{../Esonero_Big_Data/src/main/java/com/BigData/Job1_Mapper.java}
\lstinputlisting[title=Job1\_Reducer.java]{../Esonero_Big_Data/src/main/java/com/BigData/Job1_Reducer.java}

\subsection{Job 2}
\lstinputlisting[title=Job2\_Driver.java]{../Esonero_Big_Data/src/main/java/com/BigData/Job2_Driver.java}
\lstinputlisting[title=Job2\_Mapper.java]{../Esonero_Big_Data/src/main/java/com/BigData/Job2_Mapper.java}
\lstinputlisting[title=Job2\_Reducer.java]{../Esonero_Big_Data/src/main/java/com/BigData/Job2_Reducer.java}

\subsection{Job 3}
\lstinputlisting[title=Job3\_Driver.java]{../Esonero_Big_Data/src/main/java/com/BigData/Job3_Driver.java}
\lstinputlisting[title=Job3\_Mapper.java]{../Esonero_Big_Data/src/main/java/com/BigData/Job3_Mapper.java}
\lstinputlisting[title=Job3\_Reducer.java]{../Esonero_Big_Data/src/main/java/com/BigData/Job3_Reducer.java}

\end{document}